\documentclass[10pt,twocolumn]{article}

% ── Packages ──────────────────────────────────────────────────────────────
\usepackage[margin=1in]{geometry}
\usepackage{amsmath,amssymb}
\usepackage{graphicx}
\usepackage{booktabs}
\usepackage{multirow}
\usepackage{hyperref}
\usepackage{xcolor}
\usepackage{cite}
\usepackage{caption}
\usepackage{subcaption}
\usepackage{algorithm}
\usepackage{algpseudocode}
\usepackage{balance}

% ── Placeholder helpers (remove before submission) ─────────────────────────
\newcommand{\TODO}[1]{\textcolor{red}{\textbf{[TODO: #1]}}}
\newcommand{\REF}{\textcolor{blue}{[REF]}}
\newcommand{\NUM}[1]{\textcolor{teal}{[#1]}}
\newcommand{\BLANK}{\textcolor{orange}{\textbf{---}}}   % result not yet available

% ── Mask shorthand (used throughout) ──────────────────────────────────────
\newcommand{\Mfull}{\mathcal{M}_{\text{full}}}
\newcommand{\Mreg}{\mathcal{M}_{\text{reg}}}
\newcommand{\Mbone}{\mathcal{M}_{\text{bone}}}
\newcommand{\Mdice}{\mathcal{M}_{\text{Dice}}}

% ── Title ──────────────────────────────────────────────────────────────────
\title{%
  Organ-Guided Multi-Phase Abdominal CT Registration\\
  via Anatomy-Constrained Mattes Mutual Information
}

\author{%
  \TODO{Author 1}\textsuperscript{1} \quad
  \TODO{Author 2}\textsuperscript{1} \quad
  \TODO{Author 3}\textsuperscript{1,2}\\[4pt]
  \textsuperscript{1}\TODO{Institution}\\
  \textsuperscript{2}\TODO{Institution (if applicable)}\\[4pt]
  \texttt{\TODO{corresponding@email.edu}}
}

\date{}

% ══════════════════════════════════════════════════════════════════════════
\begin{document}
\maketitle

% ── Abstract ──────────────────────────────────────────────────────────────
\begin{abstract}
Accurate spatial alignment of multi-phase abdominal computed tomography
(CT) — non-contrast (NC), arterial, and venous acquisitions — is
essential for joint analysis of vascular enhancement, lesion
characterisation, and organ segmentation transfer.
The task is complicated by inter-session patient repositioning, variable
axial coverage, and the fundamental invalidity of intensity-based
similarity metrics across contrast-induced Hounsfield Unit (HU)
distributions.
We propose a fully automated, anatomy-constrained pipeline comprising
five ordered stages: (i)~preprocessing with DICOM-convention-aware
HU normalisation; (ii)~TotalSegmentator-based organ segmentation merged
into a custom four-tier label volume; (iii)~confidence-gated 3D centroid
Z-alignment; (iv)~a three-pass rigid registration scheme that initialises
from a Kabsch SVD on bone-only labels, refines via Nelder-Mead
optimisation of a combined centroid-distance and eroded-Dice loss, and
polishes with a Sobel-gradient NCC stage gated by a centroid-displacement
acceptance criterion; and (v)~B-spline deformable refinement.
Throughout, the Mattes Mutual Information (MMI) similarity objective is
evaluated exclusively within $\Mreg$, a mask of contrast-stable organs
from which bowel, contrast-enhancing structures, and peristaltic tissue
are explicitly excluded.
Applied to the VinDr multi-phase CT dataset
(\NUM{265} studies, \NUM{795} volumes, two scanner families), the
pipeline is evaluated via eroded Dice coefficient and centroid
displacement across eleven abdominal structures grouped into parenchymal,
vascular, and skeletal categories.
\TODO{Headline result: e.g., ``The proposed method achieves a median
eroded Dice of \BLANK\,$\pm$\,\BLANK across all organ groups, reducing
centroid displacement by \BLANK\,\% relative to the strongest classical
baseline, while requiring \BLANK$\times$ less wall-clock time than ANTs
SyN.''}
\end{abstract}

% ══════════════════════════════════════════════════════════════════════════
\section{Introduction}
\label{sec:intro}

Multi-phase abdominal CT is a cornerstone of hepatic oncology,
vascular surgery planning, and solid organ evaluation~\cite{forner2018}.
Each protocol comprises three acquisitions — non-contrast (NC),
arterial, and venous phases — that capture complementary tissue
properties: NC provides baseline parenchymal density, arterial phase
reveals vascular enhancement, and venous phase characterises portal
and parenchymal perfusion~\cite{easl2018}.
Joint exploitation of all three phases, whether for lesion
characterisation, longitudinal change detection, or segmentation
propagation, requires all volumes to be in accurate spatial
correspondence.

Achieving this correspondence is non-trivial.
First, patients are repositioned between acquisitions, introducing
rigid-body translations that include axial shifts up to
\NUM{80}\,mm from variable breath-hold depth and table position.
Second, contrast administration non-linearly alters the HU distribution
of vascular and parenchymal structures: the aortic lumen shifts from
$\sim$30\,HU in NC to $\sim$350\,HU in arterial phase, while liver
parenchyma enhances by 20--80\,HU in venous phase~\REF.
These shifts violate the local intensity-consistency assumptions of
standard similarity metrics such as normalised cross-correlation (NCC)
and mean squared error (MSE)~\cite{pluim2003}.
Third, bowel and hollow viscera undergo peristaltic motion and gas
redistribution between acquisitions, introducing non-rigid deformations
that corrupt any metric evaluated over the full field of view.
Prior work has addressed rigid inter-session alignment~\REF\ and
deformable abdominal registration~\cite{metz2011,klein2010}\ in
isolation, yet few methods simultaneously address all three
challenges in a single, principled pipeline.

The key gap in existing methods is the application of whole-image
similarity objectives in a setting where those objectives are
theoretically invalid.
Classical toolboxes such as elastix~\cite{klein2010} and ANTs~\cite{avants2008}
default to full-volume NCC or MSE, metrics whose validity depends on
a linear (NCC) or identity (MSE) relationship between corresponding
voxel intensities — a relationship broken wherever contrast enhancement
occurs.
Segmentation-driven approaches~\cite{xu2019,hu2018} use label overlap
as the objective, side-stepping intensity invalidity but coupling
registration accuracy directly to segmentation quality.
A principled formulation is missing: one that retains the statistical
power of intensity-based similarity but evaluates it only where
the underlying metric assumptions hold.

This paper makes the following contributions:
%
\begin{enumerate}

  \item \textbf{Anatomy-constrained MMI.}
  We propose restricting the Mattes Mutual Information objective to
  $\Mreg$, a TotalSegmentator-derived, automatically merged mask of
  contrast-stable organs, from which bowel, padding, and
  contrast-enhancing structures are explicitly excluded.
  The principled exclusion is motivated by the statistical invalidity
  of intensity similarity across contrast-induced HU shifts
  (\S\ref{sec:seg}; ablated in \S\ref{sec:metric_ablation}).

  \item \textbf{Bone-anchored Kabsch SVD initialisation.}
  Prior to intensity optimisation, we align phases via a closed-form
  Kabsch SVD~\cite{kabsch1976} computed on $\Mbone$, the bone-only
  sub-mask.
  Because bony anatomy is both phase-invariant and rigid, this
  initialisation provides a geometrically sound starting point that
  is independent of contrast-induced centroid shifts in soft tissue
  (\S\ref{sec:rigid}).

  \item \textbf{Confidence-gated 3D centroid Z-alignment.}
  We introduce a weighted-centroid Z-offset estimation stage with a
  MAD-filtered confidence score that gates fallback to a centre-of-mass
  crop when segmentation coverage is insufficient
  (\S\ref{sec:zalign}).

  \item \textbf{Organ-grouped evaluation protocol.}
  We report eroded Dice and centroid displacement per organ, grouped
  into parenchymal, vascular, and skeletal categories across
  \NUM{265} studies and two contrast phases, providing a reproducible
  benchmark for future methods on the VinDr dataset
  (\S\ref{sec:eval}).

  \item \textbf{Systematic comparison and ablation.}
  We evaluate the full pipeline against \TODO{ANTs SyN, elastix
  B-spline, and VoxelMorph} under identical evaluation conditions,
  and include an ablation over similarity metric (MMI vs.\ NCC
  vs.\ MSE, masked vs.\ whole-image) and pipeline stage contribution
  (\S\ref{sec:results}).

\end{enumerate}

% ══════════════════════════════════════════════════════════════════════════
\section{Related Work}
\label{sec:related}

\subsection{Classical Multi-Phase CT Registration}
\label{sec:related_classical}

Intensity-based deformable registration was established by Thirion's
Demons algorithm~\cite{thirion1998} and the free-form B-spline model
of Rueckert~et~al.~\cite{rueckert1999}.
The elastix toolkit~\cite{klein2010} and ANTs~\cite{avants2008}
generalised these frameworks to a broad family of similarity metrics
and transform models, becoming de facto baselines for abdominal CT
registration.
Mutual information (MI)~\cite{wells1996} and its Parzen-window
estimator~\cite{mattes2001} extended registration to multi-modal
settings by modelling the joint intensity histogram rather than
assuming linear correspondence; the statistical conditions under which
MI remains a valid registration objective are surveyed in~\cite{pluim2003}.

For multi-phase abdominal CT specifically, Metz et~al.~\cite{metz2011}
demonstrated B-spline registration with temporal groupwise optimisation
but applied the similarity metric over the full image domain.
Heinrich~et~al.~\cite{heinrich2013} introduced DEEDS, a discrete
deformable registration method that optimises over a dense displacement
field via efficient message passing, achieving strong performance
on abdominal and multi-modal data without explicit intensity assumptions.
The Learn2Reg benchmark~\cite{hering2022} provides a standardised
comparison framework for abdominal registration methods; we adopt
a compatible evaluation protocol in \S\ref{sec:eval}.

The limitation common to these classical methods is the application
of whole-image similarity objectives in the multi-phase setting, where
contrast enhancement locally violates the intensity assumptions that
justify those objectives.

\subsection{Learning-Based Abdominal Registration}
\label{sec:related_dl}

VoxelMorph~\cite{balakrishnan2019} established the amortised
learning-based registration paradigm, training a convolutional network
to predict deformation fields directly from image pairs.
TransMorph~\cite{chen2022} and LapIRN~\cite{mok2020} subsequently
demonstrated strong performance on large-deformation and multi-scale
registration.
HyperMorph~\cite{hoopes2021} further showed that registration
hyperparameters — including the weighting between similarity and
regularisation objectives — can themselves be learned, but decisions
about \emph{which regions} contribute to the similarity term remain
left to the practitioner.

A key challenge for learning-based methods in the multi-phase CT
setting is the absence of within-phase training data: most public
benchmarks provide intra-modality pairs, and cross-phase pairs require
either manual annotation or a registration method to generate
pseudo-ground-truth.
UniGradICON~\cite{tian2023} represents a step toward task-agnostic
generalisation, but its domain shift across contrast phases has
not been characterised.
The evaluation protocol introduced in \S\ref{sec:eval} is directly
applicable to these methods as future baselines.

\subsection{Anatomy-Constrained and Masked Registration}
\label{sec:related_masked}

The use of anatomical information to guide registration spans several
paradigms.
Segmentation-driven methods~\cite{xu2019,hu2018} replace the intensity
objective with a label-based loss, fully decoupling registration from
intensity validity but binding accuracy to segmentation quality.
Anatomy-constrained neural networks~\cite{oktay2017} incorporate
anatomical shape priors into the training objective, improving
plausibility of deformations near organ boundaries.

An alternative approach uses image features that are, by design,
modality-independent.
Heinrich~et~al.~\cite{heinrich2012} proposed MIND descriptors that
capture local self-similarity structure, making the registration
feature representation robust to cross-modal intensity shifts.
In the multi-phase CT context, this motivates gradient-domain
similarity (used in our Pass~2; \S\ref{sec:rigid}), since HU gradients
at organ boundaries are structurally similar across phases even when
absolute intensities differ.

Region-of-interest masking within elastix and ITK-based
pipelines~\cite{klein2010} allows the fixed-image metric domain
to be restricted to a spatial mask, but existing applications
use this feature as a spatial convenience (e.g., excluding background)
rather than as a principled exclusion of contrast-invalid tissue.
Sliding-motion regularisation~\cite{papiez2014} uses anatomical
boundaries to adapt the regulariser but evaluates similarity globally.
De~Vos~et~al.~\cite{devos2019} incorporate anatomy-aware weighting
in a learning-based framework, but the weighting is learned from
within-modality data and does not address cross-phase invalidity.

The novelty of the present work is the \emph{principled exclusion}
of contrast-sensitive and peristaltic structures from an
intensity-based objective, motivated by the statistical invalidity
of NCC and MSE across contrast-induced HU distributions.
To our knowledge, no prior work has formalised this exclusion criterion
for multi-phase abdominal CT registration, nor provided an ablation
demonstrating its impact relative to whole-image and label-based
alternatives.

% ══════════════════════════════════════════════════════════════════════════
\section{Dataset}
\label{sec:dataset}

We evaluate on the VinDr multi-phase abdominal CT dataset~\cite{vindr},
comprising \textbf{\NUM{265} studies} (\NUM{795} volumes) acquired across
two scanner families identified by DICOM UID prefixes
\texttt{1.2.840} and \texttt{1.2.392}.
Each study contains three contrast phases: \textbf{non-contrast (NC)},
\textbf{arterial}, and \textbf{venous}.
The NC phase serves as the fixed reference in all experiments; arterial
and venous phases are registered independently.

Scanner heterogeneity manifests in differing DICOM encoding conventions:
\texttt{1.2.840} series require the standard RescaleSlope/Intercept
GDCM decoding (which SimpleITK applies automatically), while
\texttt{1.2.392} series store pixels already in Hounsfield Units.
Both conventions are handled identically by the preprocessing stage
(\S\ref{sec:preproc}), which clips all volumes to
$[-1024,\,+1024]$\,HU after resampling.

\TODO{%
  Report: (1)~mean\,$\pm$\,SD slice counts per phase,
  (2)~distribution of axial offsets between NC and moving phases
  measured at the Z-alignment stage,
  (3)~number of studies successfully processed end-to-end,
  (4)~exclusion criteria and count of excluded studies.
}

% ══════════════════════════════════════════════════════════════════════════
\section{Method}
\label{sec:method}

An overview of the pipeline is shown in Fig.~\ref{fig:pipeline}.

\begin{figure}[t]
  \centering
  \TODO{Pipeline diagram: DICOM $\to$ NIfTI $\to$ TotalSeg $\to$
  Mask Merge $\to$ Crop $\to$ Z-align $\to$ Rigid (P0/P1/P2)
  $\to$ Deformable $\to$ Evaluate.
  Each box should show the output postfix and the mask it consumes.
  Distinguish $\Mreg$ (MMI domain) from $\Mdice$ (Nelder-Mead Dice
  domain) from $\Mbone$ (Kabsch SVD domain).}
  \caption{Overview of the proposed multi-phase CT registration pipeline.
  The NC phase is the fixed reference throughout. Four mask tiers derived
  from a single TotalSegmentator pass serve distinct roles at each stage.}
  \label{fig:pipeline}
\end{figure}

% ── 4.1 ───────────────────────────────────────────────────────────────────
\subsection{Preprocessing and Standardisation}
\label{sec:preproc}

Each DICOM series is read via SimpleITK's \texttt{ImageSeriesReader}
(GDCM backend) and processed as follows:
%
\begin{enumerate}
  \item \textbf{Orientation.}
  Volumes are reoriented to LPS convention, establishing a canonical
  axis ordering across scanner families.

  \item \textbf{Resampling.}
  Trilinear interpolation resamples to isotropic 1.5\,mm spacing.
  Boundary voxels are filled with $-1024$\,HU (standard CT air),
  not the scanner minimum, to prevent outside-FOV sentinel values
  (commonly $-4048$\,HU or $-3024$\,HU) from contaminating the
  mutual information histogram.

  \item \textbf{HU clipping.}
  Values are clipped to $[-1024,\,+1024]$\,HU, removing sentinel
  artefacts while retaining the full diagnostic range of abdominal
  soft tissue ($\leq\!700$\,HU for cortical bone).
  This step is applied \emph{before} saving and before TotalSegmentator
  inference.
\end{enumerate}

% ── 4.2 ───────────────────────────────────────────────────────────────────
\subsection{Organ Segmentation and Four-Tier Mask Construction}
\label{sec:seg}

TotalSegmentator v\TODO{X.X}~\cite{wasserthal2023} is applied to each
preprocessed volume, producing per-structure binary NIfTI files.
These are merged into a single multi-label volume using a fixed
project-defined integer mapping that assigns each structure a
deterministic label independent of TotalSegmentator's internal
indexing.
The mapping defines four tiers:

\medskip
\noindent\textbf{Tier~1 — Bone} (labels 21--99): spinal cord, lumbar
and lower thoracic vertebrae (L1--L5, T10--T12), sacrum, S1, hip
bones, femora, sternum, costal cartilages, and all 24 ribs.
These structures are rigid, phase-invariant, and provide the most
reliable geometric anchors for registration initialisation.

\noindent\textbf{Tier~2 — Stable soft tissue} (labels 1--19): liver,
spleen, bilateral kidneys, pancreas, gallbladder, paraspinal muscles
(autochthon bilaterally), iliopsoas bilaterally, aorta, inferior vena
cava (IVC), portal and splenic vein, and bilateral iliac arteries and
veins.
These structures maintain predictable spatial positions across contrast
phases, making their centroids and intensities (within a known HU
range) suitable for metric evaluation.

\noindent\textbf{Tier~3 — Variable soft tissue} (labels 100+): stomach,
oesophagus, urinary bladder, adrenal glands, prostate, and gluteal
muscles.
These are included in the bounding-box mask only; they are explicitly
excluded from the registration metric domain due to contrast-sensitive
HU variation or unpredictable positional change between acquisitions.

\noindent\textbf{Tier~0 — Excluded entirely}: colon, small bowel,
duodenum, and rectum.
Peristaltic motion and gas redistribution between acquisitions introduce
non-rigid deformations that are incompatible with any rigid or B-spline
model; their inclusion in the metric domain introduces conflicting
gradient signals.

\medskip
\noindent Four composite masks are derived from this label volume:

\begin{itemize}
  \item $\Mfull$: union of Tiers~1--3. Used solely to compute the
  abdominal bounding box for volume cropping; never used as a
  registration metric domain.

  \item $\Mreg$: union of Tiers~1--2 (all stable organs). Used as
  the spatial domain for MMI evaluation in all registration stages,
  and as the anchor set for Z-alignment centroid estimation.

  \item $\Mbone$: Tier~1 subset comprising spinal cord, vertebrae
  L1--L5, T10--T12, sacrum, S1, sternum, and upper ribs (left and
  right ribs~1--3). Used exclusively in the Pass~0 Kabsch SVD
  initialisation (\S\ref{sec:rigid}).

  \item $\Mdice$: curated Tier~1+2 subset comprising liver, spleen,
  bilateral kidneys, spinal cord, vertebrae L1--L5, aorta, and IVC.
  Structures with contrast-variant HU behaviour (iliac vessels) or
  high peristaltic mobility (stomach, esophagus) are excluded.
  Used exclusively in the Pass~1 Nelder-Mead Dice loss
  (\S\ref{sec:rigid}).
\end{itemize}

% ── 4.3 ───────────────────────────────────────────────────────────────────
\subsection{Confidence-Gated Z-Alignment}
\label{sec:zalign}

Inter-session axial offsets are resolved prior to intensity-based
registration by weighted centroid matching of anchor organs between
phases.

\paragraph{Centroid offset estimation.}
For anchor organ $k \in \{\text{liver, spleen, kidney}_L, \text{kidney}_R\}$
(labels 1--4 in $\Mreg$), the physical-space Z-centroid $c_k^{(Z)}$
is computed independently for NC and the moving phase.
The weighted Z-offset estimate is:
%
\begin{equation}
  \Delta Z \;=\;
  \frac{\displaystyle\sum_{k \in \mathcal{K}} w_k\,
        \bigl(c_k^{\text{NC},(Z)} - c_k^{\text{mov},(Z)}\bigr)}
       {\displaystyle\sum_{k \in \mathcal{K}} w_k}
  \label{eq:zoffset}
\end{equation}
%
where $\mathcal{K}$ is the set of anchor organs with voxel count
$\geq\!200$ in both phases, $w_\text{liver}=3.0$, and
$w_{\text{spleen}} = w_{\text{kidney}_L} = w_{\text{kidney}_R} = 2.0$.
Individual organ offsets deviating by more than $1.5\times$ the median
absolute deviation (MAD) from the median are treated as outliers and
excluded.

\paragraph{Confidence score and gating.}
A confidence score $\gamma = {\sum_{k\in\mathcal{K}} w_k} \,/\,
{\sum_k w_k^{\max}}$ reflects the fraction of anchor weight
that contributed to the estimate.
If $\gamma < 0.40$ or $|\Delta Z| > 80$\,mm, the study falls back to
a per-series organ centre-of-mass crop.
$\gamma$, $\Delta Z$, and the method applied are logged per study.

% ── 4.4 ───────────────────────────────────────────────────────────────────
\subsection{Three-Pass Organ-Masked Rigid Registration}
\label{sec:rigid}

Rigid registration comprises three sequential passes, each building on
the previous.

\paragraph{Pass~0: Kabsch SVD bone initialisation.}
Given the set of voxel coordinates $\{p_i\}$ from $\Mbone$ in the
fixed (NC) phase and $\{q_i\}$ from the corresponding labels in the
moving phase, we solve for the optimal rotation $R$ and translation
$t$ via the Kabsch algorithm~\cite{kabsch1976}:
%
\begin{equation}
  (R^*, t^*) = \operatorname*{arg\,min}_{R \in SO(3),\, t}
  \sum_{i} w_i \,\bigl\| R\,q_i + t - p_i \bigr\|^2
  \label{eq:kabsch}
\end{equation}
%
where $w_i$ is the per-label weight from $\Mbone$.
The closed-form solution factorises the weighted covariance matrix via
SVD.
Bone labels are chosen for this stage because they are rigid,
phase-invariant, and their centroids are unaffected by contrast
pooling — properties that do not hold for soft-tissue organ centroids
in multi-phase CT.
The result $T_0 = (R^*, t^*)$ is used to initialise Pass~1.

\paragraph{Pass~1: Nelder-Mead centroid + eroded-Dice refinement.}
Starting from $T_0$, a derivative-free Nelder-Mead
minimiser~\cite{nelder1965} optimises a combined loss over $\Mdice$:
%
\begin{equation}
  \mathcal{L}(T) = \lambda \cdot d_c(T) \;+\;
  \bigl(1 - \text{Dice}_\epsilon(T)\bigr)
  \label{eq:pass1loss}
\end{equation}
%
where $d_c$ is the volume-weighted centroid displacement
(Eq.~\ref{eq:centroid}), $\text{Dice}_\epsilon$ is the eroded Dice
coefficient (Eq.~\ref{eq:dice}) averaged over $\Mdice$ labels,
and $\lambda = 10.0$ balances the two terms.
Nelder-Mead is chosen over gradient-based optimisers because the
eroded-Dice term is computed on discretised binary masks and is
therefore non-differentiable.
The erosion radius per organ is set proportionally to organ size
(2.0--6.0\,mm), preventing small structures with noisy boundaries
from dominating the loss.
The result $T_1$ is used to initialise Pass~2.

\paragraph{Pass~2: Sobel-gradient NCC polish with acceptance gate.}
Pass~2 applies a phase-invariant intensity polish using SimpleITK's
\texttt{ImageRegistrationMethod} with a Sobel-gradient NCC objective:
the similarity metric is computed on the gradient-magnitude images
$|\nabla I_\text{fixed}|$ and $|\nabla I_\text{moving}|$ rather than
raw HU values.
HU gradients at organ boundaries are structurally similar across
contrast phases even when absolute intensities differ, providing a
phase-invariant intensity signal.
Multi-resolution proceeds at shrink factors $[3, 2, 1]$ with Gaussian
smoothing $[1.5, 0.5, 0.0]$\,mm, initialised from $T_1$.
The result $T_2$ is accepted only when centroid displacement under
$T_2$ is lower than under $T_1$; otherwise $T_1$ is retained.
This acceptance gate prevents the polish stage from corrupting
well-converged rigid alignments.

Segmentation masks are propagated to registered space via
nearest-neighbour resampling; volumes are resampled with Lanczos
windowed sinc interpolation.

% ── 4.5 ───────────────────────────────────────────────────────────────────
\subsection{B-Spline Deformable Registration}
\label{sec:deformable}

Deformable refinement uses a free-form B-spline
transform~\cite{rueckert1999} at 25\,mm control-point grid spacing,
applied after the rigid stage.
The similarity metric is MMI within $\Mreg$ (fixed mask only; no
moving mask is applied, as simultaneous masking introduces conflicting
constraints before masks are aligned).
Optimisation uses L-BFGS-B with multi-resolution $[4, 2, 1]$ and
\NUM{X} histogram bins and \NUM{Y}\% spatial sampling.

\TODO{If grid-spacing ablation (15/25/40\,mm) was run, add one sentence
referencing \S\ref{sec:ablation}; otherwise remove.}

% ══════════════════════════════════════════════════════════════════════════
\section{Evaluation Protocol}
\label{sec:eval}

\subsection{Metrics}

All metrics are computed between the NC fixed phase and each registered
moving phase (arterial, venous), summarised as median\,$\pm$\,IQR.
Median and IQR are used throughout as they are robust to the subset
of studies where segmentation fails for one phase.

\textbf{Eroded Dice coefficient.}
Both fixed and moving organ masks are independently eroded by
$r = 6$\,mm in physical space prior to Dice computation:
%
\begin{equation}
  \text{Dice}_\epsilon(A,B) =
  \frac{2\,|A_\epsilon \cap B_\epsilon|}{|A_\epsilon| + |B_\epsilon|}
  \label{eq:dice}
\end{equation}
%
Erosion removes the boundary voxels where TotalSegmentator uncertainty
concentrates, ensuring the metric reflects volumetric alignment quality
rather than delineation accuracy.
Organ pairs in which either eroded mask contains fewer than 500 voxels
are excluded from that study's statistics.

\textbf{Centroid displacement (mm).}
Euclidean distance between physical-space centroids of the full
(non-eroded) organ masks:
%
\begin{equation}
  d_c = \bigl\|\bar{\mathbf{x}}^{(A)} - \bar{\mathbf{x}}^{(B)}\bigr\|_2
  \label{eq:centroid}
\end{equation}
%
where $\bar{\mathbf{x}}^{(\cdot)}$ is the voxel-count-weighted centroid
in mm.
This metric is insensitive to boundary noise and sensitive to gross
misalignment, making it complementary to eroded Dice.

\textbf{NCC exclusion.}
Normalised cross-correlation assumes a linear relationship between
corresponding voxel intensities, which is violated across contrast
phases~\cite{pluim2003}.
NCC values are noise rather than signal before intensity normalisation
and are excluded from all reported results.

\subsection{Organ Groups}

The eleven evaluated structures are grouped as follows:

\begin{center}
\small
\begin{tabular}{ll}
\toprule
\textbf{Group} & \textbf{Structures} \\
\midrule
Parenchymal & Liver, Spleen, Kidney (L/R), Pancreas, Gallbladder \\
Vascular    & Aorta, IVC \\
Skeletal    & Vertebra L1, L2, L3 \\
\bottomrule
\end{tabular}
\end{center}

Vertebral bodies are included as contrast-independent internal
sanity references: correct rigid registration should yield near-perfect
skeletal alignment independent of the similarity metric choice.

\subsection{Evaluated Conditions}

\textbf{Pipeline ablation} (same framework, varying stages):

\begin{center}
\small
\begin{tabular}{ll}
\toprule
\textbf{Stage} & \textbf{Description} \\
\midrule
Baseline   & Organ CoM crop, no registration \\
Z-Aligned  & Z-centroid alignment only \\
Rigid      & Z-align + three-pass rigid \\
Deformable & Z-align + rigid + B-spline \\
\bottomrule
\end{tabular}
\end{center}

\textbf{Metric ablation} (deformable stage, fixed organ mask):
MMI (proposed) vs.\ NCC vs.\ MSE, each evaluated whole-image and
within $\Mreg$. All other hyperparameters held fixed.

\textbf{Method comparison} (deformable stage):
\TODO{ANTs SyN~\cite{avants2008}, elastix B-spline~\cite{klein2010},
DEEDS~\cite{heinrich2013}, VoxelMorph~\cite{balakrishnan2019}}.
Wall-clock time per study (GPU/CPU) reported alongside accuracy.

% ══════════════════════════════════════════════════════════════════════════
\section{Experiments and Results}
\label{sec:results}

\subsection{Implementation Details}
\label{sec:impl}

The pipeline is implemented in Python~\NUM{3.X} using
SimpleITK~\NUM{v2.X}~\cite{lowekamp2013} and
TotalSegmentator~\NUM{vX.X}~\cite{wasserthal2023}.
Segmentation runs on a single \TODO{GPU model} (\TODO{X}\,GB VRAM).
Registration (all three passes and B-spline) runs on CPU.
\TODO{Report mean\,$\pm$\,SD wall-clock time per study per stage:
segmentation / Z-align / rigid (P0 + P1 + P2) / deformable.}

\subsection{Pipeline Stage Ablation}
\label{sec:ablation}

Table~\ref{tab:ablation} reports eroded Dice and centroid displacement
across all organ groups and phases at each pipeline stage.

\begin{table*}[t]
\centering
\caption{%
  Registration quality at each pipeline stage.
  Eroded Dice (higher is better) and centroid displacement in mm
  (lower is better), reported as median\,$\pm$\,IQR across
  \NUM{N} studies per phase. NC is the fixed reference. \BLANK{} =
  result pending.
}
\label{tab:ablation}
\setlength{\tabcolsep}{5pt}
\begin{tabular}{ll cccc cccc cccc}
\toprule
& & \multicolumn{4}{c}{\textbf{Parenchymal}} &
    \multicolumn{4}{c}{\textbf{Vascular}} &
    \multicolumn{4}{c}{\textbf{Skeletal}} \\
\cmidrule(lr){3-6}\cmidrule(lr){7-10}\cmidrule(lr){11-14}
\textbf{Stage} & \textbf{Ph.} &
  \multicolumn{2}{c}{Dice $\uparrow$} & \multicolumn{2}{c}{CD$\downarrow$} &
  \multicolumn{2}{c}{Dice $\uparrow$} & \multicolumn{2}{c}{CD$\downarrow$} &
  \multicolumn{2}{c}{Dice $\uparrow$} & \multicolumn{2}{c}{CD$\downarrow$} \\
\midrule
Baseline    & A & \BLANK & \BLANK & \BLANK & \BLANK & \BLANK & \BLANK & \BLANK & \BLANK & \BLANK & \BLANK & \BLANK & \BLANK \\
            & V & \BLANK & \BLANK & \BLANK & \BLANK & \BLANK & \BLANK & \BLANK & \BLANK & \BLANK & \BLANK & \BLANK & \BLANK \\
\midrule
Z-Aligned   & A & \BLANK & \BLANK & \BLANK & \BLANK & \BLANK & \BLANK & \BLANK & \BLANK & \BLANK & \BLANK & \BLANK & \BLANK \\
            & V & \BLANK & \BLANK & \BLANK & \BLANK & \BLANK & \BLANK & \BLANK & \BLANK & \BLANK & \BLANK & \BLANK & \BLANK \\
\midrule
Rigid       & A & \BLANK & \BLANK & \BLANK & \BLANK & \BLANK & \BLANK & \BLANK & \BLANK & \BLANK & \BLANK & \BLANK & \BLANK \\
            & V & \BLANK & \BLANK & \BLANK & \BLANK & \BLANK & \BLANK & \BLANK & \BLANK & \BLANK & \BLANK & \BLANK & \BLANK \\
\midrule
Deformable  & A & \BLANK & \BLANK & \BLANK & \BLANK & \BLANK & \BLANK & \BLANK & \BLANK & \BLANK & \BLANK & \BLANK & \BLANK \\
            & V & \BLANK & \BLANK & \BLANK & \BLANK & \BLANK & \BLANK & \BLANK & \BLANK & \BLANK & \BLANK & \BLANK & \BLANK \\
\bottomrule
\end{tabular}
\end{table*}

\TODO{%
  Analysis paragraph: Which stage provides the largest gain?
  Does Z-alignment alone close most of the gap, or is intensity
  registration necessary? Do vascular structures improve less than
  parenchymal (expected due to smaller volume)? Do vertebrae confirm
  rigid correctness?
}

\subsection{Similarity Metric Ablation}
\label{sec:metric_ablation}

Table~\ref{tab:metric} isolates the contribution of the anatomy
constraint from the choice of similarity metric by fixing the pipeline
stage (deformable) and varying both the metric and its evaluation
domain.

\begin{table}[t]
\centering
\caption{%
  Effect of similarity metric and masking on deformable registration
  quality. All conditions use identical pipeline stages and
  hyperparameters. Median eroded Dice / centroid displacement (mm).
  \BLANK{} = result pending.
}
\label{tab:metric}
\begin{tabular}{lcccc}
\toprule
\textbf{Condition} & \multicolumn{2}{c}{\textbf{Arterial}} &
                     \multicolumn{2}{c}{\textbf{Venous}} \\
\cmidrule(lr){2-3}\cmidrule(lr){4-5}
 & Dice $\uparrow$ & CD $\downarrow$ & Dice $\uparrow$ & CD $\downarrow$ \\
\midrule
MSE (whole image)     & \BLANK & \BLANK & \BLANK & \BLANK \\
NCC (whole image)     & \BLANK & \BLANK & \BLANK & \BLANK \\
MSE (organ mask)      & \BLANK & \BLANK & \BLANK & \BLANK \\
NCC (organ mask)      & \BLANK & \BLANK & \BLANK & \BLANK \\
MMI (whole image)     & \BLANK & \BLANK & \BLANK & \BLANK \\
\textbf{MMI (organ mask)} & \textbf{\BLANK} & \textbf{\BLANK} & \textbf{\BLANK} & \textbf{\BLANK} \\
\bottomrule
\end{tabular}
\end{table}

\TODO{%
  Analysis: Show (a) masking consistently improves over whole-image
  for any metric, and (b) MMI is best within the masked setting.
  This is the core empirical argument of the paper.
}

\subsection{Comparison with Existing Methods}
\label{sec:comparison}

\begin{table*}[t]
\centering
\caption{%
  Comparison against existing registration methods at the
  deformable/non-rigid stage on \NUM{265} studies.
  CD = centroid displacement (mm). Time = mean wall-clock
  seconds per study (CPU unless noted). \textbf{Bold} = best per column.
  \BLANK{} = result pending.
}
\label{tab:comparison}
\begin{tabular}{l ccc ccc ccc c}
\toprule
& \multicolumn{3}{c}{\textbf{Parenchymal}} &
  \multicolumn{3}{c}{\textbf{Vascular}} &
  \multicolumn{3}{c}{\textbf{Skeletal}} &
  \textbf{Time (s)} \\
\cmidrule(lr){2-4}\cmidrule(lr){5-7}\cmidrule(lr){8-10}
\textbf{Method} &
  Dice $\uparrow$ & CD $\downarrow$ && Dice $\uparrow$ & CD $\downarrow$ &&
  Dice $\uparrow$ & CD $\downarrow$ && \\
\midrule
ANTs SyN~\cite{avants2008}
  & \BLANK & \BLANK && \BLANK & \BLANK && \BLANK & \BLANK && \BLANK \\
elastix B-spline~\cite{klein2010}
  & \BLANK & \BLANK && \BLANK & \BLANK && \BLANK & \BLANK && \BLANK \\
DEEDS~\cite{heinrich2013}
  & \BLANK & \BLANK && \BLANK & \BLANK && \BLANK & \BLANK && \BLANK \\
VoxelMorph~\cite{balakrishnan2019}
  & \BLANK & \BLANK && \BLANK & \BLANK && \BLANK & \BLANK && \BLANK \\
\midrule
\textbf{Proposed}
  & \textbf{\BLANK} & \textbf{\BLANK} && \textbf{\BLANK} & \textbf{\BLANK} &&
    \textbf{\BLANK} & \textbf{\BLANK} && \BLANK \\
\bottomrule
\end{tabular}
\end{table*}

\TODO{%
  Analysis: Where does the proposed method outperform baselines most
  clearly? Any cases where a baseline wins? Discuss the time--accuracy
  trade-off.
}

\subsection{Qualitative Results}
\label{sec:qual}

\begin{figure*}[t]
  \centering
  \TODO{%
    Figure 2: Axial overlay panels for a representative study.
    Show NC (fixed) and arterial (moving) at four stages:
    Baseline / Z-Aligned / Rigid / Deformable.
    Use checkerboard or edge overlay. Annotate liver, aorta, L2.
    Show one easy and one hard (large Z-offset) case.
  }
  \caption{Qualitative registration results for a representative study
  (axial mid-liver slice). NC fixed phase shown in green, arterial
  moving phase in magenta. Misalignment appears as colour fringing;
  correct alignment produces a grey overlay.}
  \label{fig:qual}
\end{figure*}

\TODO{One paragraph referencing Fig.~\ref{fig:qual}. Highlight
progression from Baseline to Deformable. Note residual misalignment
in difficult structures (pancreas, gallbladder) and explain why.}

% ══════════════════════════════════════════════════════════════════════════
\section{Discussion}
\label{sec:discussion}

\TODO{%
  D1 — Anatomy-constrained MMI.
  Argue from Table~\ref{tab:metric}: masking matters more than metric
  choice (i.e., MMI masked $>$ MMI whole-image $\geq$ NCC masked),
  or vice versa depending on results. Explain mechanistically:
  in multi-phase CT, the joint intensity histogram between NC and
  arterial liver contains two distinct modes (enhancing vs.
  non-enhancing parenchyma). Whole-image MMI marginalises over this
  heterogeneity; organ-masked MMI evaluates it where tissue identity
  is known.
}

\TODO{%
  D2 — Kabsch SVD initialisation.
  Argue that Pass~0 provides a geometrically correct warm start that
  prevents Pass~1 from converging to local optima in the Dice/centroid
  loss surface. If the Pass~0 ablation is available (rigid with and
  without Kabsch init), quantify the gain here.
}

\TODO{%
  D3 — Organ-group observations.
  Parenchymal: liver and spleen show the largest absolute improvement.
  Pancreas and gallbladder: small volume, contrast-variable, expect
  higher variance.
  Vascular: aorta is rigid and contrast-bright in arterial; expect
  good alignment. IVC is compressible — may show higher centroid
  displacement.
  Skeletal: any residual centroid error here indicates a registration
  failure, not an anatomy challenge.
}

\TODO{%
  D4 — Limitations.
  (1) No deep learning component — iterative inference.
  Report total study time; discuss suitability for retrospective
  cohort analysis vs. real-time.
  (2) Pipeline depends on TotalSegmentator quality; segmentation
  failures propagate downstream (characterise failure rate).
  (3) Evaluation uses pseudo ground truth (TotalSegmentator masks);
  no manually annotated landmarks.
  (4) Extension to other datasets (CHAOS, LiTS, multi-centre trials)
  requires verifying DICOM encoding conventions.
}

% ══════════════════════════════════════════════════════════════════════════
\section{Conclusion}
\label{sec:conclusion}

We presented a multi-phase abdominal CT registration pipeline whose
central design principle is the confinement of intensity similarity
to an automatically derived mask of contrast-stable organs.
By evaluating Mattes Mutual Information exclusively within $\Mreg$
— a mask from which bowel, contrast-enhancing vascular tissue, and
peristaltic structures are explicitly excluded — the pipeline avoids
the corrupting influence of phase-dependent HU variation on the
registration objective.
A bone-anchored Kabsch SVD initialisation and a derivative-free
Nelder-Mead refinement stage, together with a Sobel-gradient NCC
acceptance gate, provide geometrically robust convergence across
scanner types and large axial offsets.
\TODO{Quantitative summary: e.g., ``On \NUM{265} VinDr studies, the
proposed method achieves \BLANK\,\% higher median eroded Dice and
\BLANK\,mm lower centroid displacement than the strongest classical
baseline.''}
The four-tier mask construction and eroded-Dice evaluation protocol
provide a reproducible benchmark for future registration methods on
multi-phase abdominal CT.

% ══════════════════════════════════════════════════════════════════════════
\bibliographystyle{IEEEtran}
% \bibliography{references}  % Uncomment and use .bib file before submission

\begin{thebibliography}{99}

% ── Classical registration ─────────────────────────────────────────────
\bibitem{rueckert1999}
D.~Rueckert \emph{et~al.},
``Nonrigid registration using free-form deformations: Application to
breast MR images,''
\emph{IEEE Trans.\ Med.\ Imag.}, vol.~18, no.~8, pp.~712--721, 1999.

\bibitem{thirion1998}
J.-P.~Thirion,
``Image matching as a diffusion process: An analogy with Maxwell's
demons,''
\emph{Med.\ Image Anal.}, vol.~2, no.~3, pp.~243--260, 1998.

\bibitem{avants2008}
B.~B.~Avants \emph{et~al.},
``Symmetric diffeomorphic image registration with cross-correlation:
Evaluating automated labeling of elderly and neurodegenerative brain,''
\emph{Med.\ Image Anal.}, vol.~12, no.~1, pp.~26--41, 2008.

\bibitem{klein2010}
S.~Klein \emph{et~al.},
``Elastix: A toolbox for intensity-based medical image registration,''
\emph{IEEE Trans.\ Med.\ Imag.}, vol.~29, no.~1, pp.~196--205, 2010.

\bibitem{mattes2001}
D.~Mattes \emph{et~al.},
``Nonrigid multimodality image registration,''
in \emph{Proc.\ SPIE Med.\ Imag.}, vol.~4322, 2001, pp.~1609--1620.

\bibitem{wells1996}
W.~M.~Wells~III \emph{et~al.},
``Multi-modal volume registration by maximization of mutual information,''
\emph{Med.\ Image Anal.}, vol.~1, no.~1, pp.~35--51, 1996.

\bibitem{pluim2003}
J.~P.~W.~Pluim, J.~B.~A.~Maintz, and M.~A.~Viergever,
``Mutual-information based registration of medical images: A survey,''
\emph{IEEE Trans.\ Med.\ Imag.}, vol.~22, no.~8, pp.~986--1004, 2003.

\bibitem{studholme1999}
C.~Studholme, D.~L.~G.~Hill, and D.~J.~Hawkes,
``An overlap invariant entropy measure of 3D medical image alignment,''
\emph{Pattern Recognition}, vol.~32, no.~1, pp.~71--86, 1999.

\bibitem{metz2011}
C.~T.~Metz \emph{et~al.},
``Nonrigid registration of dynamic medical imaging data using nD+t
B-splines and a groupwise optimization approach,''
\emph{Med.\ Image Anal.}, vol.~15, no.~2, pp.~238--249, 2011.

\bibitem{heinrich2013}
M.~P.~Heinrich \emph{et~al.},
``DEEDS: MRF-based nonlinear image registration using efficient discrete
optimization,''
\emph{Int.\ J.\ Comput.\ Vis.}, vol.~103, no.~3, pp.~384--400, 2013.

\bibitem{hering2022}
A.~Hering \emph{et~al.},
``Learn2Reg: Comprehensive multi-task medical image registration challenge,
dataset and evaluation,''
\emph{IEEE Trans.\ Med.\ Imag.}, vol.~41, no.~6, pp.~1458--1476, 2022.

\bibitem{murphy2011}
K.~Murphy \emph{et~al.},
``Semi-automatic construction of reference standards for evaluation of
image registration,''
\emph{Med.\ Image Anal.}, vol.~15, no.~1, pp.~71--84, 2011.

% ── Learning-based registration ───────────────────────────────────────
\bibitem{balakrishnan2019}
G.~Balakrishnan \emph{et~al.},
``VoxelMorph: A learning framework for deformable medical image
registration,''
\emph{IEEE Trans.\ Med.\ Imag.}, vol.~38, no.~8, pp.~1788--1800, 2019.

\bibitem{chen2022}
J.~Chen \emph{et~al.},
``TransMorph: Transformer for unsupervised medical image registration,''
\emph{Med.\ Image Anal.}, vol.~82, p.~102615, 2022.

\bibitem{mok2020}
T.~C.~W.~Mok and A.~Chung,
``Large deformation diffeomorphic image registration with Laplacian
pyramid networks,''
in \emph{Proc.\ MICCAI}, 2020, pp.~211--221.

\bibitem{hoopes2021}
A.~Hoopes \emph{et~al.},
``HyperMorph: Amortised hyperparameter learning for image registration,''
in \emph{Proc.\ IPMI}, 2021, pp.~3--17.

\bibitem{tian2023}
H.~Tian \emph{et~al.},
``UniGradICON: A foundation model for medical image registration,''
\emph{arXiv:2403.05780}, 2024.
% TODO: Check for published venue; arXiv as placeholder.

\bibitem{zhao2019}
Z.~Zhao \emph{et~al.},
``Recursive cascaded networks for unsupervised medical image registration,''
in \emph{Proc.\ ICCV}, 2019, pp.~10600--10610.

\bibitem{xu2019}
Z.~Xu and M.~Niethammer,
``DeepAtlas: Joint semi-supervised learning of image registration and
segmentation,''
in \emph{Proc.\ MICCAI}, 2019, pp.~420--429.

% ── Anatomy-constrained / masked registration ─────────────────────────
\bibitem{hu2018}
Y.~Hu \emph{et~al.},
``Weakly-supervised convolutional neural networks for multimodal image
registration,''
\emph{Med.\ Image Anal.}, vol.~49, pp.~1--13, 2018.

\bibitem{oktay2017}
O.~Oktay \emph{et~al.},
``Anatomically constrained neural networks (ACNN): Application to
cardiac image enhancement and segmentation,''
\emph{IEEE Trans.\ Med.\ Imag.}, vol.~37, no.~2, pp.~384--395, 2018.
% Note: Published in IEEE TMI 2018; arXiv submitted 2017.

\bibitem{heinrich2012}
M.~P.~Heinrich \emph{et~al.},
``MIND: Modality independent neighbourhood descriptor for multi-modal
deformable registration,''
\emph{Med.\ Image Anal.}, vol.~16, no.~7, pp.~1423--1435, 2012.

\bibitem{papiez2014}
B.~W.~Papiez \emph{et~al.},
``An implicit sliding-motion preserving regularisation via bilateral
filtering for deformable image registration,''
\emph{Med.\ Image Anal.}, vol.~18, no.~8, pp.~1323--1335, 2014.

\bibitem{devos2019}
B.~D.~de~Vos \emph{et~al.},
``A deep learning framework for unsupervised affine and deformable
image registration,''
\emph{Med.\ Image Anal.}, vol.~52, pp.~128--143, 2019.

% ── Tools and infrastructure ──────────────────────────────────────────
\bibitem{wasserthal2023}
J.~Wasserthal \emph{et~al.},
``TotalSegmentator: Robust segmentation of 104 anatomical structures
in CT images,''
\emph{Radiol.:\ Artif.\ Intell.}, vol.~5, no.~5, 2023.

\bibitem{lowekamp2013}
B.~C.~Lowekamp \emph{et~al.},
``The design of SimpleITK,''
\emph{Front.\ Neuroinform.}, vol.~7, p.~45, 2013.

\bibitem{kabsch1976}
W.~Kabsch,
``A solution for the best rotation to relate two sets of vectors,''
\emph{Acta Crystallogr.\ A}, vol.~32, no.~5, pp.~922--923, 1976.

\bibitem{nelder1965}
J.~A.~Nelder and R.~Mead,
``A simplex method for function minimization,''
\emph{Comput.\ J.}, vol.~7, no.~4, pp.~308--313, 1965.

% ── Clinical motivation ───────────────────────────────────────────────
\bibitem{forner2018}
A.~Forner, M.~Reig, and J.~Bruix,
``Hepatocellular carcinoma,''
\emph{Lancet}, vol.~391, no.~10127, pp.~1301--1314, 2018.

\bibitem{easl2018}
European Association for the Study of the Liver,
``EASL clinical practice guidelines: Management of hepatocellular
carcinoma,''
\emph{J.\ Hepatol.}, vol.~69, no.~1, pp.~182--236, 2018.

% ── Dataset ───────────────────────────────────────────────────────────
\bibitem{vindr}
\TODO{VinDr multi-phase CT dataset citation. Search PhysioNet or
  authors' publications for the correct reference. Placeholder.}

\end{thebibliography}

\end{document}